\begin{exercise}[VI.11.E01 --- Local-ring criterion]
Prove that if every element $a$ of a ring $R$ satisfies: either $a$ or $1-a$ is a unit, then $R$ is local.
\end{exercise}

\begin{exercise}[VI.11.E02 --- Units in the affine-line local ring]
Describe explicitly the units of $k[t]_{(t-a)}$ and identify its maximal ideal.
\end{exercise}

\begin{exercise}[VI.11.E03 --- Prime localization]
For $\mathfrak p\in\Spec A$, prove directly that $a/s\in A_{\mathfrak p}$ is a unit iff $a\notin\mathfrak p$.
\end{exercise}

\begin{exercise}[VI.11.E04 --- Local spectrum]
Show that primes of $A_{\mathfrak p}$ correspond exactly to primes $\mathfrak q\subseteq\mathfrak p$ of $A$.
\end{exercise}

\begin{exercise}[VI.11.E05 --- Residue field formula]
Compute $\kappa(\mathfrak p)$ from $A/\mathfrak p$ and prove $\kappa(\mathfrak p)\cong\Frac(A/\mathfrak p)$.
\end{exercise}

\begin{exercise}[VI.11.E06 --- Affine line]
Compute the local rings and residue fields at $(0)$ and $(t-a)$ in $\Spec k[t]$.
\end{exercise}

\begin{exercise}[VI.11.E07 --- Real quadratic point]
For $\mathfrak m=(t^2+1)\subset\RR[t]$, compute $\RR[t]_{\mathfrak m}$ and its residue field.
\end{exercise}

\begin{exercise}[VI.11.E08 --- Arithmetic line]
Compute $\ZZ_{(p)}$, its units, maximal ideal, and residue field.
\end{exercise}

\begin{exercise}[VI.11.E09 --- Finite-field closed point]
If $f\in\FF_q[t]$ is irreducible of degree $d$, compute the residue field at $(f)$.
\end{exercise}

\begin{exercise}[VI.11.E10 --- Parabola generic point]
Compute the local ring and residue field of $k[x,y]$ at the prime $(y-x^2)$.
\end{exercise}

\begin{exercise}[VI.11.E11 --- Crossing at origin]
Let $A=k[x,y]/(xy)$.  Show that $A_{(x,y)}$ has zero divisors and residue field $k$.
\end{exercise}

\begin{exercise}[VI.11.E12 --- Crossing away from origin]
In the same ring, compute the local ring at $(x-a,y)$ for $a\neq0$.
\end{exercise}

\begin{exercise}[VI.11.E13 --- Minimal prime]
For $A=k[x,y]/(xy)$ compute $A_{(x)}$ and $A_{(y)}$.
\end{exercise}

\begin{exercise}[VI.11.E14 --- Dual numbers]
Show $k[\varepsilon]/(\varepsilon^2)$ is local; compute its maximal ideal, units, and residue field.
\end{exercise}

\begin{exercise}[VI.11.E15 --- Finite local/nonlocal comparison]
Determine which of $\ZZ/4\ZZ$, $\ZZ/6\ZZ$, and $\ZZ/8\ZZ$ are local and compute residue fields of their maximal ideals.
\end{exercise}

\begin{exercise}[VI.11.E16 --- Shrinking invariance]
If $f\notin\mathfrak p$, construct explicitly the isomorphism $(A_f)_{\mathfrak pA_f}\cong A_{\mathfrak p}$.
\end{exercise}

\begin{exercise}[VI.11.E17 --- Quotient-localization]
If $I\subseteq\mathfrak p$, prove $(A/I)_{\mathfrak p/I}\cong A_{\mathfrak p}/IA_{\mathfrak p}$.
\end{exercise}

\begin{exercise}[VI.11.E18 --- Residue field after quotient]
Under the hypotheses of the preceding exercise, prove the residue fields at $\mathfrak p$ and $\mathfrak p/I$ are naturally isomorphic.
\end{exercise}

\begin{exercise}[VI.11.E19 --- Zero germ criterion]
Prove $a/1=0$ in $A_{\mathfrak p}$ iff $sa=0$ for some $s\notin\mathfrak p$.
\end{exercise}

\begin{exercise}[VI.11.E20 --- Local map]
Given $\varphi:A\to B$ and $\mathfrak q\in\Spec B$, with $\mathfrak p=\varphi^{-1}(\mathfrak q)$, construct $A_{\mathfrak p}\to B_{\mathfrak q}$ and prove it is local.
\end{exercise}

\begin{exercise}[VI.11.E21 --- Residue-field map]
Continue the preceding exercise and construct the induced injective field map $\kappa(\mathfrak p)\to\kappa(\mathfrak q)$.
\end{exercise}

\begin{exercise}[VI.11.E22 --- Product ring]
Compute the localizations of $k\times k$ at its two prime ideals.
\end{exercise}

\begin{exercise}[VI.11.E23 --- Mixed arithmetic point]
For $A=\ZZ[x]$ and $\mathfrak m=(2,x-1)$, compute $A/\mathfrak m$ and $\kappa(\mathfrak m)$, and describe which elements are units in $A_{\mathfrak m}$.
\end{exercise}

\begin{exercise}[VI.11.E24 --- Specialization chain]
For $(0)\subset(x)\subset(x,y)$ in $k[x,y]$, construct the maps $A_{(x,y)}\to A_{(x)}\to A_{(0)}$ and explain why the direction reverses the prime inclusions.
\end{exercise}
